% Authority: exact-scope leaf 8; full PDF page 87; printed p.70.
% U:p070.u01 | paragraph-start
\par\indent
We shall thus have:
\[
p=t^2+2u^2=\varphi^2+\psi^2,
\]
and, after transposing:
\[
t^2-\psi^2=(t+\psi)(t-\psi)=\varphi^2-2u^2.
\]
% U:p070.u02 | paragraph-start
\par\indent
Since $\varphi$ is odd, the greatest common divisor of $\varphi$ and $u$ will be odd;
designate it by $m$ and put $\varphi=m\varphi'$, $u=mu'$. Substitution of these
values in the preceding equation changes it into:
\[
(t+\psi)(t-\psi)=m^2(\varphi'^2-2u'^2).
\]
% U:p070.u03 | continuation
The odd number $t+\psi$ is thus composed of two factors $E$ and $K$, the first
of which divides $m^2$, the second $\varphi'^2-2u'^2$; and the same is true of
$t-\psi$, whose factors we denote by $F$ and $L$, $L$ being allowed to be negative.
We therefore have the equations:
\[
\begin{aligned}
t+\psi&=EK, &t-\psi&=FL,\\
m^2&=EF, &\varphi'^2-2u'^2&=KL.
\end{aligned}
\]
% U:p070.u04 | continuation
It is easy to verify that the numbers $E$ and $F$ are relatively prime. Indeed,
let $\delta$ be a prime divisor of $E$, and suppose that $\delta$ also divides $F$.
The number $\delta$ would be a common divisor of $t+\psi$ and $t-\psi$, and consequently
would divide the number $t$, which is the half-sum of the preceding two numbers.
On the other hand, from the fact that $\delta$ is a prime divisor of $E$, it follows
successively, having regard to the equations $EF=m^2$, $u=mu'$, that $m^2$, $m$,
and $u$ are multiples of $\delta$. The numbers $t$ and $u$ would therefore have the
common divisor $\delta$, and $p=t^2+2u^2$ would not be a prime number.
% U:p070.u05 | paragraph-start
\par\indent
The product of the numbers $E$ and $F$, which are relatively prime, being a
square, each of them is also a square. Let $E=e^2$; then $t+\psi=e^2K$. The odd
square $e^2$ is of the form $8n+1$, and $K$, as an odd divisor of
$\varphi'^2+2u'^2$ (where $\varphi'$ and $u'$ are relatively prime), is of one of
these forms: $8n+1$, $8n+7$. Thus the number $t+\psi$ will itself be of one of
the forms $8n+1$, $8n+7$. The number $\psi$, which we know to be divisible by $4$,
will be of the form $8n$ or of the form $8n+4$. It follows from the preceding
that, in the first case, $t$ will be of one of these two forms: $8n+1$, $8n+7$,
and in the second of one of these: $8n+3$, $8n+5$. Comparing this result with
the preceding theorem, one has this other theorem:
% U:p070.u06 | paragraph-start
\par\indent
“$p$ designating a prime number $8n+1$, and having written $p=\varphi^2+\psi^2$
(where $\psi$ is supposed divisible by $4$), $\pm2$ will or will not be a biquadratic
residue with respect to $p$ according as $\psi$ is of the form $8n$ or of the form $8n+4$.”
